% Generated draft: requires mathematical and source review.
\documentclass[11pt]{article}
\usepackage[margin=25mm]{geometry}
\usepackage{fontspec}
\setmainfont{Latin Modern Roman}
\usepackage{amsmath,amssymb,mathtools,unicode-math}
\setmathfont{Latin Modern Math}
\usepackage{enumitem,needspace,xurl,hyperref}
\setlength{\parindent}{0pt}
\setlength{\parskip}{5pt}
\setlength{\emergencystretch}{3em}
\newcommand{\N}{\mathbb{N}}
\newcommand{\Z}{\mathbb{Z}}
\newcommand{\Q}{\mathbb{Q}}
\newcommand{\R}{\mathbb{R}}
\newcommand{\C}{\mathbb{C}}
\newcommand{\F}{\mathbb{F}}
\newcommand{\A}{\mathbb{A}}
\newcommand{\PP}{\mathbb{P}}
\newcommand{\E}{\mathbb{E}}
\newcommand{\eps}{\varepsilon}
\newcommand{\dd}{\,\mathrm{d}}
\begin{document}
\section*{UnsolvedMath ID: 30001996}
\textbf{Source problem code:} catalogue entry
\subsection{Definitions and mathematical statement}
The supplied record asks an open-ended question about a parameterized ensemble of controlled systems; it does not specify its equations, parameter space, state spaces, admissible controls, or meaning of controllability. The following is a conditional formal schema, not a uniquely determined formulation of the source problem.

Let \(\Lambda\) be a parameter set and, for each \(\lambda\in\Lambda\), let \(X_\lambda\) be a state space. Once an evolution law and admissibility conditions have been specified, write \(\Phi_\lambda(T,x;u)\) for the state reached from \(x\in X_\lambda\) at time \(T>0\) under a control \(u\). Denote by \(\mathcal U_\lambda(T)\) the controls admissible for the individual system and by \(\mathcal U_{\mathrm{ens}}(T)\) the common controls admissible simultaneously for the ensemble. These sets, and the existence domains of the evolution maps, are unspecified in the supplied record.

Under the interpretation of individual controllability as exact controllability between arbitrary states in finite, unrestricted time, the premise would be
\[
\forall\lambda\in\Lambda\;\forall x,y\in X_\lambda\;\exists T>0\;\exists u\in\mathcal U_\lambda(T):
\quad \Phi_\lambda(T,x;u)=y.
\]
Here the time and control may depend on \(\lambda,x,y\). To formulate exact joint controllability, one must additionally specify a class \(\mathcal P\subseteq\prod_{\lambda\in\Lambda}X_\lambda\) of admissible state profiles. One possible conclusion is
\[
\forall a,b\in\mathcal P\;\exists T>0\;\exists u\in\mathcal U_{\mathrm{ens}}(T)\;\forall\lambda\in\Lambda:
\quad \Phi_\lambda\bigl(T,a(\lambda);u\bigr)=b(\lambda).
\]
Thus a single control and a single terminal time must work for every parameter, although they may depend on the pair of profiles.

\textbf{Problem, subject to the unresolved interpretation above.} Determine additional hypotheses on the ensemble, its parameter dependence, and its admissible controls under which controllability of every individual system implies joint controllability of the ensemble.

The source does not determine whether either displayed exact-controllability interpretation is intended. If joint controllability instead means approximate controllability, the terminal profile must be required to approach the target in a specified topology, with the control and terminal time quantified for every prescribed neighborhood of that target. Neither that topology nor the admissible profile class is supplied. The source also leaves unresolved whether time is prescribed, whether controllability is from a distinguished initial state, and whether the ensemble is finite or infinite. Consequently, no unique fully specified implication can be extracted without additional source information.
\subsection{Short English statement}
When each member of an ensemble can be controlled separately, determine what additional assumptions allow one shared control to steer the ensemble to its target profile. The supplied record leaves the system class and the exact or approximate meaning of joint controllability unspecified, so the quantified formulation is conditional and requires human mathematical review.
\subsection{Sources}
\begin{itemize}
\item Local registry: supplied UnsolvedMath source record, permanent ID \(30001996\), source problem code \texttt{catalogue entry}, title \emph{Ensemble Controllability from Individual Controllability}, in the Oberwolfach Reports Open Problems collection. The record gives proposed year 2012 and names no individual proposer. Its original question is: ``What can be said about controllability of the joint system when every individual system in the ensemble is controllable?''
\item Supplied report citation: \emph{Control Theory: Mathematical Perspectives on Complex Networked Systems} (2012), Oberwolfach Reports, pp.\ 661--732, DOI \texttt{10.4171/owr/2012/12}. \url{https://doi.org/10.4171/owr/2012/12}. The record describes Oberwolfach Reports as documenting mathematical workshops and their open-problem discussions; it does not identify a specific page or named contributor for this question.
\item Supplied literature source, labeled ``Oberwolfach Report 12 2012'': \url{https://doi.org/10.4171/owr/2012/12}.
\item Supplied literature source, labeled ``Ensemble Control on Lie Groups'': \url{https://arxiv.org/abs/2008.03243}.
\item Literature assessment checked at \texttt{2026-08-21}: the supplied dated curation labels the problem \texttt{partially\_solved} and states that general smooth ensembles remain broader than the known theory, while the individual-to-ensemble implication is proved for important bilinear ensembles on semisimple Lie groups and related special classes. This is supplied dated curation derived from the OWR record, not an independently verified literature assessment; the time-sensitive status and cited claims require independent verification.
\item The registry supplies statement status \texttt{corrected\_verified} and reports a manual semantic, typographic, and LaTeX review of its cleaned wording. That supplied editorial assessment does not resolve the missing mathematical specifications or constitute independent verification here.
\end{itemize}
\end{document}
